\pdfvariable trailerid {[<C3492026090300000000000000000000><C3492026090300000000000000000000>]}
\documentclass[10pt]{article}
\usepackage[margin=0.72in]{geometry}
\usepackage{amsmath,amsthm,unicode-math,microtype,xurl,booktabs}
\providecommand{\CRevisionRound}{2}
\newtheorem{theorem}{Theorem}
\newtheorem{lemma}[theorem]{Lemma}
\newcommand{\RR}{\mathbb R}
\newcommand{\TT}{\mathbb T}
\setlength{\emergencystretch}{3em}
\title{Uhlenbeck Integrals, a Rational Lax Matrix, and\\
Compact Liouville Fibers for the Neumann Sphere}
\author{Route-A source-local certificate HCS-C349}
\date{3 September 2026}

\begin{document}
\maketitle
\begin{abstract}
For the anisotropic Neumann oscillator on the round two-sphere we prove
global completeness and derive all three Uhlenbeck integrals with their two
linear relations.
\ifnum\CRevisionRound>0
A rational $2\times2$ Lax matrix generates the integrals as determinant
residues.  Direct Dirac-bracket involution and compactness close the regular
Liouville two-tori and their necessary-and-sufficient rational-frequency
return criterion.
\fi
\ifnum\CRevisionRound>1
The six axial equilibria, invariant coordinate circles, repeated-spectrum
$SO(2)$ faces, isotropic great circles, and natural compact quantization are
closed separately.  Exact finite receipts test conventions but do not
replace the continuum proof.
\fi
\end{abstract}

\section{Frozen constrained system and theorem}
Let $A=\operatorname{diag}(a_1,a_2,a_3)$ with $a_1<a_2<a_3$.  On
\[
 T^*S^2=\{(x,p)\in\RR^3\times\RR^3:|x|^2=1,
                  \ x\mathbin{\cdot}p=0\}
\]
take
\begin{equation}\label{eq:N}
 H=\frac12(|p|^2+x^TAx),\qquad
 \dot x=p,\qquad
 \dot p=-Ax+\alpha x,\quad \alpha=x^TAx-|p|^2.
\end{equation}
Put $L_{ij}=x_ip_j-x_jp_i$ and
\begin{equation}\label{eq:F}
 F_i=x_i^2+\sum_{j\ne i}\frac{L_{ij}^2}{a_i-a_j}.
\end{equation}

\begin{theorem}\label{thm:main}
The flow~\eqref{eq:N} is complete.  The $F_i$ are conserved, pairwise
Dirac--Poisson commuting, and satisfy
\[
 \sum_iF_i=1,\qquad \sum_i a_iF_i=2H.
\]
\ifnum\CRevisionRound>0
They are the residues of the rational Lax determinant in
Section~\ref{sec:lax}.  Every connected regular common fiber of two
independent integrals is a Liouville two-torus; the physical flow on it is
periodic precisely when its frequency vector has a common period.
\fi
\ifnum\CRevisionRound>1
The axial, coordinate, repeated-spectrum, and isotropic clauses in
Section~\ref{sec:boundary} exhaust the declared parameter boundaries.  The
source Hamiltonian has the natural self-adjoint compact-resolvent
quantization stated there.
\fi
\end{theorem}

\section{Completeness and the Uhlenbeck identities}
The two constraint derivatives are
\[
 \frac d{dt}|x|^2=2x\mathbin{\cdot}p,\qquad
 \frac d{dt}(x\mathbin{\cdot}p)=|p|^2-x^TAx+\alpha=0
\]
on $T^*S^2$.  Also $\dot H=0$.  On $H=E$ the identity
$|p|^2=2E-x^TAx$ bounds $p$; the orbit remains in a compact subset, and the
smooth vector field is complete.

The angular momenta satisfy
\begin{equation}\label{eq:Ldot}
 \dot L_{ij}=(a_i-a_j)x_ix_j.
\end{equation}
Hence
\[
 \dot F_i=2x_i\left(p_i+\sum_{j\ne i}L_{ij}x_j\right)=0,
\]
because the parenthesis equals
$p_i+x_i(x\mathbin{\cdot}p)-p_i|x|^2$.  Pairing $(i,j)$ and $(j,i)$ proves
$\sum F_i=1$.  The weighted pairing and the Gram identity give
\[
 \sum_i a_iF_i=x^TAx+\sum_{i<j}L_{ij}^2
 =x^TAx+|p|^2=2H.
\]
This is the \emph{compact Uhlenbeck owner} added in revision round zero.

\ifnum\CRevisionRound>0
\section{Resolvent Lax generator}\label{sec:lax}
For $\lambda\notin\{a_1,a_2,a_3\}$ let
\[
 R_\lambda=(\lambda I-A)^{-1},\quad
 U=x^TR_\lambda x,\quad V=x^TR_\lambda p,\quad
 W=1+p^TR_\lambda p
\]
and define
\begin{equation}\label{eq:lax}
 \mathcal L=\begin{pmatrix}V&U\\-W&-V\end{pmatrix},\qquad
 \mathcal M=\begin{pmatrix}0&1\\\alpha-\lambda&0\end{pmatrix}.
\end{equation}
Using $a_i/(\lambda-a_i)=\lambda/(\lambda-a_i)-1$ in
\eqref{eq:N} gives
\[
 \dot U=2V,\qquad
 \dot V=W+(\alpha-\lambda)U,
 \qquad \dot W=2(\alpha-\lambda)V.
\]
These are exactly $\dot{\mathcal L}=[\mathcal L,\mathcal M]$.  Moreover,
\begin{equation}\label{eq:residue}
 \det\mathcal L=UW-V^2
 =\sum_{i=1}^3\frac{F_i}{\lambda-a_i}.
\end{equation}
Indeed, the residue at $a_i$ is~\eqref{eq:F}; both rational functions
vanish at infinity and have no other poles.

The constrained bracket is
\[
 \{x_i,x_j\}_D=0,\quad
 \{x_i,p_j\}_D=\delta_{ij}-x_ix_j,\quad
 \{p_i,p_j\}_D=x_jp_i-x_ip_j.
\]
Differentiating~\eqref{eq:F}, substituting these three coordinate brackets,
and pairing the denominators $a_i-a_j$ and $a_j-a_i$ gives
\begin{equation}\label{eq:commute}
 \{F_i,F_j\}_D=0\qquad(1\le i,j\le3).
\end{equation}
The full finite-sum cancellation is recorded in the proof package and is
independently expanded by the symbolic lane.

The two linear relations leave two generically independent integrals.  A
regular common fiber is closed in a compact energy shell.  Liouville--Arnold
therefore identifies each connected component with $\TT^2$ and gives
$\theta(t)=\theta(0)+t\omega$.  Return occurs exactly when
\begin{equation}\label{eq:return}
 T\omega\in2\pi\mathbb Z^2\quad\hbox{for some }T>0.
\end{equation}
For two nonzero components this is $\omega_1/\omega_2\in\mathbb Q$.
Continuously selected resonant tori are not an isolated primitive-orbit
ledger.  Equations~\eqref{eq:lax}--\eqref{eq:return} are the
\emph{resolvent torus owner} of revision round one.
\fi

\ifnum\CRevisionRound>1
\section{All declared degenerations and quantization}\label{sec:boundary}
At $x=\pm e_i,p=0$, a tangent component $q_j$, $j\ne i$, obeys
\[
 \ddot q_j=-(a_j-a_i)q_j.
\]
Thus the two copies over $e_1,e_2,e_3$ are respectively
elliptic--elliptic, saddle--center, and saddle--saddle.  Every face
$x_i=p_i=0$ is invariant and reduces to
\[
 H_i=\frac12\dot\theta^2+\frac12
     (a_j\cos^2\theta+a_k\sin^2\theta),
 \qquad \{i,j,k\}=\{1,2,3\}.
\]

If $a_1=a_2$, the individual fractions $F_1,F_2$ are not used.
Equation~\eqref{eq:Ldot} instead preserves the $SO(2)$ Noether momentum
$J_{12}=L_{12}$, while $F_3$ remains nonsingular and $SO(2)$ invariant.  Thus,
writing $a_1=a_2=a\ne b=a_3$,
\begin{equation}\label{eq:double}
 \{J_{12},F_3\}_D=0,\qquad
 2H=a+J_{12}^2+(b-a)F_3.
\end{equation}
This commuting pair is independent on a nonempty open set.  Indeed, at
$x=(3/5,0,4/5)$ and $p=(0,t,0)$, its differential determinant on the tangent
variations $((-4/5,0,3/5),0)$ and $(0,(0,1,0))$ is
\[
 -\frac{8\bigl(9(b-a)+25t^2\bigr)}{125(b-a)},
\]
which is nonzero after avoiding the sole possible exceptional value of
$t>0$.  The other double faces follow by permutation.  If $A=aI$, then
$\dot p=-|p|^2x$: $p=0$ is the
equilibrium sphere and $p\ne0$ is a great circle of least period
$2\pi/|p|$.

Fix a positive Planck parameter $\hbar>0$.  The canonical source quantization is
\begin{equation}\label{eq:quantum}
 \widehat H=-\frac{\hbar^2}{2}\Delta_{S^2}+\frac12x^TAx,
 \qquad D(\widehat H)=H^2(S^2)\subset L^2(S^2).
\end{equation}
The round Laplacian is self-adjoint with compact resolvent, and the real
smooth potential is bounded.  The bounded-perturbation theorem preserves
both properties.  The excluded value $\hbar=0$ would instead require the
maximal $L^2(S^2)$ domain for the bounded multiplication operator.  No closed
anisotropic spectrum is asserted.

\section{Evidence, sources, and Route-A boundary}
The exact ledger has 60 rational tangent-state rows, 120 rational resolvent
probes, 30 axial rows, 30 coordinate-face rows, and explicit axisymmetric
and isotropic boundaries.  A producer-independent checker recomputes every
fraction.  A separate symbolic lane expands the Poisson brackets, resolvent
identity, and Lax equation.  These receipts test implementation conventions;
the preceding proof handles the continuum.

The nearest workspace owners are C186 Euler top, C244 spherical pendulum,
C313 round-sphere geodesics, and C344 resonant-triad elliptic dynamics.  None
has this holonomic anisotropic sphere potential and its Uhlenbeck--resolvent
generator.

The conservative tuple is
\[
 (\mathrm{A0\_FAIL},\mathrm{A1\_WEAK},\mathrm{A2\_FAIL},
  \mathrm{A3\_FAIL},\mathrm{A4\_NATURAL\_QUANTIZATION}).
\]
There is no rational-prime carrier, logarithmic-prime clock, target Euler
factor, root number, automorphy, target divisor or functional equation,
target-zero match, or Hilbert--Polya operator.  Route A is rejected and
Route B is false.  This is the \emph{boundary and route firewall} of the
final revision.
\fi

\section*{Source boundary}
Neumann owns the original mechanical problem; Moser owns the modern
quadrics/spectral lineage; Knoerrer owns the geodesic correspondence cited
here.  This convention-locked reconstruction makes no priority claim.

\begin{thebibliography}{9}
\bibitem{Neumann}
C. Neumann, ``De problemate quodam mechanico, quod ad primam integralium
ultraellipticorum classem revocatur,'' \emph{J. Reine Angew. Math.} 56
(1859), 46--63, DOI \url{10.1515/crll.1859.56.46}.
\bibitem{Moser}
J. Moser, ``Geometry of Quadrics and Spectral Theory,'' in
\emph{The Chern Symposium 1979}, Springer (1980), 147--188,
DOI \url{10.1007/978-1-4613-8109-9_7}.
\bibitem{Knoerrer}
H. Knoerrer, ``Geodesics on quadrics and a mechanical problem of C. Neumann,''
\emph{J. Reine Angew. Math.} 334 (1982), 69--78,
DOI \url{10.1515/crll.1982.334.69}.
\end{thebibliography}
\end{document}
